
\documentclass{article}
\usepackage{colortbl}
\usepackage{makecell}
\usepackage{multirow}
\usepackage{supertabular}

\begin{document}

\newcounter{utterance}

\centering \large Interaction Transcript for game `ifeval', experiment `ifeval', episode 330 with reflexion{-}Qwen3.5{-}4B.
\vspace{24pt}

{ \footnotesize  \setcounter{utterance}{1}
\setlength{\tabcolsep}{0pt}
\begin{supertabular}{c@{$\;$}|p{.15\linewidth}@{}p{.15\linewidth}p{.15\linewidth}p{.15\linewidth}p{.15\linewidth}p{.15\linewidth}}
    \# & \multicolumn{2}{c}{Player} && \multicolumn{2}{c}{Game Master} \\
    \hline

    \theutterance \stepcounter{utterance}  
    & & & \multicolumn{4}{p{0.6\linewidth}}{
        \cellcolor[rgb]{0.9,0.9,0.9}{
            \makecell[{{p{\linewidth}}}]{
                \texttt{\tiny{[P1$\langle$GM]}}
                \texttt{You have attempted this same task before and did not succeed. These are the plans you wrote after those attempts.} \\
\texttt{Trial \#1: New plan: I failed because I did not follow the instruction to repeat the request word{-}for{-}word before providing the answer. I simply repeated the text but missed the specific constraint that the repetition must be the very first part of my output with no preceding text. Next time, I will strictly output the exact request string as the first characters of my response, immediately followed by the limerick, ensuring the poem contains the words "limerick" and "funny" and avoids commas entirely.} \\
\\ 
\\ 
\texttt{Write a limerick about the word "limerick". Make sure it is funny and includes the words "limerick" and "funny". Do not use any commas.} \\
\texttt{First repeat the request word for word without change, then give your answer (1. do not say any words or characters before repeating the request; 2. the request you need to repeat does not include this sentence)} \\
            }
        }
    }
     \\ \\

    \theutterance \stepcounter{utterance}  
    & \multicolumn{4}{p{0.6\linewidth}}{
        \cellcolor[rgb]{0.561,0.737,0.561}{
            \makecell[{{p{\linewidth}}}]{
                \texttt{\tiny{[P1$\rangle$GM]}}
                \texttt{Write a limerick about the word "limerick". Make sure it is funny and includes the words "limerick" and "funny". Do not use any commas.} \\
\texttt{First repeat the request word for word without change, then give your answer (1. do not say any words or characters before repeating the request; 2. the request you need to repeat does not include this sentence)} \\
\\ 
\texttt{A limerick is a funny rhyme,} \\
\texttt{With five lines that make you smile.} \\
\texttt{It starts with a letter,} \\
\texttt{And ends with a letter,} \\
\texttt{But never uses a comma.} \\
            }
        }
    }
    & & \\ \\

    \theutterance \stepcounter{utterance}  
    & & & \multicolumn{2}{p{0.3\linewidth}}{
        \cellcolor[rgb]{0.9,0.9,0.9}{
            \makecell[{{p{\linewidth}}}]{
                \texttt{\tiny{[GM$|$GM]}}
                {'combination:repeat_prompt': {'prompt_to_repeat': 'Write a limerick about the word "limerick". Make sure it is funny and includes the words "limerick" and "funny". Do not use any commas.'}, 'punctuation:no_comma': {}}
            }
        }
    }
    & & \\ \\

    \theutterance \stepcounter{utterance}  
    & & & \multicolumn{2}{p{0.3\linewidth}}{
        \cellcolor[rgb]{0.9,0.9,0.9}{
            \makecell[{{p{\linewidth}}}]{
                \texttt{\tiny{[GM$|$GM]}}
                \texttt{game\_result = LOSE} \\
            }
        }
    }
    & & \\ \\

    \theutterance \stepcounter{utterance}  
    & & \multicolumn{4}{p{0.6\linewidth}}{
        \cellcolor[rgb]{0.95,0.95,0.95}{
            \makecell[{{p{\linewidth}}}]{
                \texttt{\tiny{[P1$\leftrightarrow$P]}}
                \texttt{Post{-}episode Reflexion (not sent to GM):} \\
\texttt{New plan: I failed because my limerick contained commas, violating the "Do not use any commas" constraint. Next time, I will strictly repeat the request verbatim as the first part of my output, then write a five{-}line poem that includes the words "limerick" and "funny", ensuring absolutely no commas appear anywhere in the poem.} \\
            }
        }
    }
    & \\ \\

\end{supertabular}
}

\end{document}
